\documentclass[11pt]{article}
\usepackage[a4paper,margin=2cm]{geometry}
\usepackage{amsmath,amssymb}
\usepackage[hidelinks]{hyperref}
\allowdisplaybreaks
\title{Qpu.General.Meaning.Config}
\author{Tsvetan Rouschev}
\date{@uuidna/qpu v1.1.1, CC-BY-NC-ND-4.0}
\begin{document}
\maketitle
\noindent\url{https://github.com/uuidna/qpu/blob/v1.1.1/lean/Qpu/General/Meaning/Config.lean}

\paragraph{\texttt{agentConfirmRules\_eq\_cons}}
The list whose first element is `head`, where `tail` is the rest of the list. Usually written `head :: tail`.   Conventions for notations in identifiers:   * The recommended spelling of `::` in identifiers is `cons`.   * The recommended spelling of `[a]` in identifiers is `singleton`.
\[ \mathrm{Config.agentConfirmRules} = \mathrm{Config.agentConfirmRules.headD}\ \mathrm{default} :: \mathrm{Config.agentConfirmRules.tail} \]
\noindent axioms: none; proof: \url{https://github.com/uuidna/qpu/blob/v1.1.1/lean/Qpu/General/Meaning/Config.lean\#L59}

\paragraph{\texttt{agentEvents\_eq\_cons}}
The list whose first element is `head`, where `tail` is the rest of the list. Usually written `head :: tail`.   Conventions for notations in identifiers:   * The recommended spelling of `::` in identifiers is `cons`.   * The recommended spelling of `[a]` in identifiers is `singleton`.
\[ \mathrm{Config.agentEvents} = \mathrm{Config.agentEvents.headD}\ \mathrm{default} :: \mathrm{Config.agentEvents.tail} \]
\noindent axioms: none; proof: \url{https://github.com/uuidna/qpu/blob/v1.1.1/lean/Qpu/General/Meaning/Config.lean\#L26}

\paragraph{\texttt{agentSpawnTools\_eq\_cons}}
The list whose first element is `head`, where `tail` is the rest of the list. Usually written `head :: tail`.   Conventions for notations in identifiers:   * The recommended spelling of `::` in identifiers is `cons`.   * The recommended spelling of `[a]` in identifiers is `singleton`.
\[ \mathrm{Config.agentSpawnTools} = \mathrm{Config.agentSpawnTools.headD}\ \mathrm{default} :: \mathrm{Config.agentSpawnTools.tail} \]
\noindent axioms: none; proof: \url{https://github.com/uuidna/qpu/blob/v1.1.1/lean/Qpu/General/Meaning/Config.lean\#L48}

\paragraph{\texttt{agentWriteTools\_eq\_cons}}
The list whose first element is `head`, where `tail` is the rest of the list. Usually written `head :: tail`.   Conventions for notations in identifiers:   * The recommended spelling of `::` in identifiers is `cons`.   * The recommended spelling of `[a]` in identifiers is `singleton`.
\[ \mathrm{Config.agentWriteTools} = \mathrm{Config.agentWriteTools.headD}\ \mathrm{default} :: \mathrm{Config.agentWriteTools.tail} \]
\noindent axioms: none; proof: \url{https://github.com/uuidna/qpu/blob/v1.1.1/lean/Qpu/General/Meaning/Config.lean\#L37}

\paragraph{\texttt{tenants\_eq\_cons}}
The list whose first element is `head`, where `tail` is the rest of the list. Usually written `head :: tail`.   Conventions for notations in identifiers:   * The recommended spelling of `::` in identifiers is `cons`.   * The recommended spelling of `[a]` in identifiers is `singleton`.
\[ \mathrm{Config.tenants} = \mathrm{Config.tenants.headD}\ \mathrm{default} :: \mathrm{Config.tenants.tail} \]
\noindent axioms: none; proof: \url{https://github.com/uuidna/qpu/blob/v1.1.1/lean/Qpu/General/Meaning/Config.lean\#L15}
\end{document}
